\begin{abstract}
\noindent
We test whether the normal modes of the U.S. equity correlation matrix
exhibit predictable mean reversion. In the overdamped limit of a
coupled-oscillator market, the covariance matrix is the inverse of the
stiffness matrix, so the lowest-variance eigenportfolios should relax
fastest. Using \DataAssets{} S\&P~500 stocks, monthly walk-forward
re-estimation and out-of-sample data from 2013 to 2024, we find that the
modes the Marchenko--Pastur filter certifies as signal do not revert
(Sharpe \SRSignalLevel), while modes deep in the noise bulk do (Sharpe
up to \SRMiddleLevel, $t=\TMiddleLevel$), with significantly negative
daily autocorrelation. Supervised models using eigenvalue separation,
half-life, autocorrelation, eigenvector stability, volatility,
portfolio concentration and regime features predict which modes will
revert with out-of-sample AUC of 0.51--0.53: significant, and better
than random selection, but not better than holding every mode. A
two-state hidden Markov model reveals a strong regime asymmetry:
reversion profits arise in high-volatility, high-correlation states and
reverse in calm ones. A walk-forward strategy combining mode
predictions with regime-based exposure reaches a net Sharpe ratio of
\SRMLOverB{} at 2~bp costs. That result depends on same-close execution
and on an overlay whose form was chosen after observing the regime
asymmetry. Overall, the normal-mode picture locates mean reversion
correctly in the spectrum, but after costs, and once the specifications
we tried are taken into account, it provides little tradable alpha.
\end{abstract}
